%% Overleaf			
%% Software Manual and Technical Document Template	
%% 									
%% This provides an example of a software manual created in Overleaf.

\documentclass{../ol-softwaremanual}

% Packages used in this example
\usepackage{graphicx}  % for including images
\usepackage{microtype} % for typographical enhancements
\usepackage{minted}    % for code listings
\usepackage{amsmath}   % for equations and mathematics
\setminted{style=friendly,fontsize=\small}
\renewcommand{\listoflistingscaption}{List of Code Listings}
\usepackage{hyperref}  % for hyperlinks
\usepackage[a4paper,top=4.2cm,bottom=4.2cm,left=3.5cm,right=3.5cm]{geometry} % for setting page size and margins

\usepackage[english, greek]{babel}

\usepackage{subfig}

\usepackage{incgraph,tikz}

\usepackage{filemod}





\usepackage{rotating}


% Custom macros used in this example document
\newcommand{\doclink}[2]{\href{#1}{#2}\footnote{\url{#1}}}
\newcommand{\cs}[1]{\texttt{\textbackslash #1}}

\begin{document}
	
	
	\begin{titlepage}
		
		
		% Frontmatter data; appears on title page
		\title{\en Team Plan \\}
		\version{0.1}
		\softwarelogo{\includegraphics[scale=0.4]{../CarBazaar_logo.png}}		
		
	\end{titlepage}
	
	
	\maketitle
	
	\newpage
	
	\center{\textbf{Μέλη Ομάδας}}
	
	\vspace{20pt}
	
	
	
	\begin{table}[htbp!]
		
		\begin{tabular}{llll}
			Μεμελετζόγλου Χαρίλαος & 1069364 & \en st1069364@ceid.upatras.gr & 4o Έτος   \\ 
			\\ Λέκκας Γεώργιος      &      1067430    &   \en st1067430@ceid.upatras.gr & 4o Έτος  \\
			\\ Γιαννουλάκης Ανδρέας        &   1067387       & \en st1067387@ceid.upatras.gr & 4o Έτος           \\
			\\ Κανελλόπουλος Ιωακείμ        &  1070914        &    \en st1070914@ceid.upatras.gr & 4o Έτος        \\ 
		\end{tabular}
	\end{table}
	
	\center{\textbf{Υπεύθυνοι Παρόντος Τεχνικού Κειμένου}}
	
	\vspace{20pt}
	
	\begin{table}[htbp!]
		\begin{tabular}{ll}
			Μεμελετζόγλου Χαρίλαος & \en Editor \\
			\\ Λέκκας Γεώργιος      &   \en  Peer Reviewer \\
			\\ Γιαννουλάκης Ανδρέας & \en Contributor \\
			\\ Καννελόπουλος Ιωακείμ & \en Contributor \\ 
		\end{tabular}
	\end{table}
	
	
	\vspace{20pt}
	
	\center{\textbf{Εργαλεία που χρησιμοποιήθηκαν}}
	
	\vspace{20pt}
	\flushleft
	Χρησιμοποιήθηκε το \en \doclink{https://www.overleaf.com/}{Overleaf} \gr και το \en \doclink{https://www.texstudio.org/}{TexStudio} \gr για την συγγραφή του \LaTeX\ κώδικα. \break
	
	Για την δημιουργία του λογότυπου, χρησιμοποιήθηκε το εργαλείο \en \doclink{https://www.adobe.com/express/create/logo}{Adobe Express} . \gr \break
	
	Για την ανάπτυξη του έργου χρησιμοποιείται η γλώσσα προγραμματισμού \en Python \gr και για την συγγραφή του κώδικα, το \en IDE \doclink{https://www.jetbrains.com/pycharm/}{PyCharm} \gr και το \en \doclink{https://code.visualstudio.com/}{VSCode} \gr .         \\ \break
	
	Για την δημιουργία του \en Gantt Chart, \gr χρησιμοποιήθηκε το \en \doclink{https://instagantt.com/}{Instagantt}  \gr και για την δημιουργία του \en Pert Chart \gr, το \en \doclink{https://www.smartdraw.com/}{Smartdraw} \gr. \break
	
	Για την παρακολούθηση των \en To-Do tasks \gr, των ολοκληρωμένων,κλπ, δηλ. για την δημιουργία ενός \en Scrum Board \gr, θα χρησιμοποιηθεί το \en site \doclink{https://trello.com/}{Trello} \gr. \break 
	
	Για την επικοινωνία των μελών της ομάδας, χρησιμοποιείται το \en \doclink{ https://www.discord.com/}{Discord} \gr . \linebreak 
	
	
	Για τον διαμοιρασμό και την παρακολούθηση της διαδικασίας ανάπτυξης, χρησιμοποιείται το \en \doclink{https://github.com/st1069364/CarBazaar}{Github} \gr.
	
	
	
	\newpage
	
	\center{\textbf{Διαδικασία Διαχείρισης Έργου}}
	\flushleft
	Προκειμένου να οργανώσουμε την εργασία της ομάδας μας, αποφασίσαμε να δουλέψουμε ακολουθώντας την μέθοδο \en \textit{SCRUM} \gr , που ανήκει στην κατηγορία των \en Agile development frameworks. \gr
	
	Συγκεκριμένα, αποφασίσαμε να εργαστούμε χρησιμοποιώντας \en \textit{Sprints} \gr διάρκειας δύο εβδομάδων, όση και η προθεσμία για το κάθε παραδοτέο. \\ 
	
	Ρόλοι κατά \en Scrum \gr : \newline
	
	
	\flushleft
	\begin{tabular}{ll}
		\en \textbf{Project Owner : }  & \gr \hspace{5mm}  Μεμελετζόγλου Χαρίλαος \\
		\\ \en \textbf{Scrum Master : } &  \gr \hspace{5mm} Λέκκας Γεώργιος \\
		
		\\ \en \textbf{Development Team : } & \begin{tabular}[t]{ll}
			&  \gr  Μεμελετζόγλου Χαρίλαος  \\
			& \gr     Λέκκας Γεώργιος \\
			& \gr     Γιαννουλάκης Ανδρέας \\
			& \gr     Κανελλόπουλος Ιωακείμ \\
		\end{tabular} 
	\end{tabular} \linebreak
	
	\vspace{20pt}
	
	Στην αρχή κάθε \en Sprint cycle \gr , διεξάγεται το \en Sprint Planning meeting \gr , στο οποίο καθορίζονται οι στόχοι της ομάδας για το \en Sprint \gr που πρόκειται να ξεκινήσει. Συγκεκριμένα, ο \en Project Owner \gr , επιλέγει τα \en tasks \gr με τα οποία θα ασχοληθεί η ομάδα, από το \en Project Backlog \gr . \\
	
	\vspace{5pt}
	
	Στόχος σε κάθε \en Sprint Cycle \gr , είναι η ανάπτυξη νέων \en features \gr και η παράδοση ενός \en working product \gr στους πελάτες. \\ 
	
	\vspace{5pt}
	
	Στην συνέχεια, η ομάδα ανάπτυξης, ξεκινά να εργάζεται στα καθορισμένα \en tasks \gr .
	Για να είναι συντονισμένη η διαδικασία ανάπτυξης των απαραίτητων Τεχνικών Κειμένων για το κάθε παραδοτέο, διεξάγονται 2 με 3 εβδομαδιαία \en Scrum meetings \gr . Στις συναντήσεις, αυτές, γίνεται έλεγχος προόδου, κάθε μέλος της ομάδας εξηγεί ποια \en tasks \gr έχει ολοκληρώσει και ποιά θα αναλάβει στην συνέχεια, αλλά και συζητά τυχόν εμπόδια που αντιμετωπίζει.
	Η επικοινωνία μεταξύ των μελών της ομάδας γίνεται μέσω \en Discord \gr, αλλα και μέσω δια ζώσης συναντήσεων. \\
	
	\vspace{5pt}
	
	Στο τέλος κάθε \en Sprint Cycle \gr , διεξάγεται το \en Sprint Retrospective meeting \gr , στο οποίο η ομάδα συζητά σχετικά με τις δυσκολίες και τα λάθη που έγιναν στο \en Sprint \gr που μόλις ολοκληρώθηκε, αλλά και πιθανούς τρόπους για να αντιμετωπιστούν αντίστοιχα φαινόμενα σε μελλοντικά \en Sprints \gr . Επίσης, γίνεται συζήτηση επί των μεθόδων που είχαν ως αποτέλεσμα την ομαλότερη διεκπεραίωση των \en tasks \gr του \en Sprint Backlog \gr .
	
	\newpage
	
		\center{\textbf{Χρονοπρογραμματισμός Έργου}}
		
		\flushleft
		
		Παρακάτω παρατίθεται ο αρχικός χρονοπρογραμματισμός του έργου. Ως έναρξη του \en project \gr έχει οριστεί η 1η Μαρτίου και ως λήξη η 12η Ιουνίου. Αξίζει να σημειωθεί πως η ημερομηνία λήξης του έργου, αποτελεί απλώς μια εκτίμηση, καθώς ενδέχεται να αλλάξει από τους διδάσκοντες. Επίσης, η χρονική διάρκεια του κάθε \en task \gr μέσα στο κάθε \en Sprint Cycle \gr, ενδέχεται να αλλάξει κατά την διάρκεια του κάθε \en Sprint \gr . Τέλος, μιας και σε κάθε παραδοτέο υπάρχουν και προαιρετικά Τεχνικά Κείμενα, ενδέχεται να προστεθούν στο μέλλον.
	
	
		\begin{figure}[htbp!]
			
			\includegraphics[width=\textwidth,height=\textheight,keepaspectratio]{img/gantt\_chart.png}
			\caption{Αρχικός Χρονοπρογραμματισμός Έργου}
			
			
		\end{figure}
	
		Ως τυπικά υποέργα (\en tasks \gr), έχουμε τα Τεχνικά Κείμενα που πρέπει να παραδοθούν σε κάθε ένα από τα παραδοτέα. Τα \en tasks \gr, φαίνονται στο διάγραμμα, αν και μπορεί η λίστα με τα \en tasks \gr κάθε παραδοτέου, να μεταβληθεί στο μέλλον, λόγω προσθήκης προαιρετικών υποέργων. 
		
		Στο τέλος κάθε παραδοτέου (και άρα κάθε \en Sprint Cycle \gr) έχουμε ένα \en Milestone \gr. Συγκεκριμένα, έχουμε :
		
		\begin{enumerate}
			\item Ανάλυση Απαιτήσεων Εφαρμογής (\en Requirements Engineering\gr)
			\item Σχεδιασμός Γενικής Αρχιτεκτονικής και Λειτουργιών Συστήματος
			\item  Επέκταση Αρχιτεκτονικής και \en Minimum Viable Product \gr.
			\item \en Testing \gr Υλοποίησης
			\item Τελικές αλλαγές στο σύστημα και παράδοση στον πελάτη
		\end{enumerate}
		\newpage
		
		
		\center{\textbf{\en Pert Chart \gr Έργου}}
		\flushleft
		Παρακάτω υπάρχει το αρχικό \en Pert Chart \gr, όπως αυτό προκύπτει από το αντίστοιχο \en Gantt Chart\gr. Μια αρχική εκτίμηση για το \en Critical Path \gr φαίνεται με κόκκινες ακμές.
		
		\begin{figure}[htbp!]
			
			\includegraphics[width=\textwidth, height=\textheight, keepaspectratio ]{img/PertChart\_for\_TeamPlan.jpg}
						\caption{Αρχικό \en Pert Chart \gr Έργου}
		\end{figure}
	


	
\end{document}
